% Problem record op_3596f8c955c66d94. This TeX file is derived from
% database/problems_json/op_3596f8c955c66d94.json by scripts/sync-tex.mjs; edit the JSON record
% and rerun the script rather than editing this file.
\section{Polynomial-time quantum algorithm for Learning With Errors}
\label{sec:op_3596f8c955c66d94}

\paragraph{Problem.}
Does the Learning With Errors problem in its standard worst-case-hard
parameter regime admit a polynomial-time quantum algorithm?

Let $n$ be the dimension, let $q=q(n)\geq 2$ be an integer modulus, and let
$\alpha=\alpha(n)\in(0,1)$ be a noise rate.  A search-LWE instance is
generated by choosing a secret $s\in\mathbb Z_q^n$ uniformly at random and
providing polynomially many independent samples
\begin{equation}
  (a_i,b_i)\in\mathbb Z_q^n\times\mathbb Z_q,
  \qquad
  b_i=\langle a_i,s\rangle+e_i\pmod q,
  \label{eq:3596-lwe-samples}
\end{equation}
where each $a_i$ is uniform in $\mathbb Z_q^n$, and each error $e_i$ is
sampled independently from a discrete Gaussian of width $\alpha q$.  For
$x\in\mathbb Z$,
\begin{equation}
  \Pr[e_i=x]=\frac{\exp\left(-\pi x^2/(\alpha q)^2\right)}
    {\sum_{z\in\mathbb Z}\exp\left(-\pi z^2/(\alpha q)^2\right)} .
  \label{eq:3596-lwe-error}
\end{equation}
The task is to recover the secret $s$ from samples satisfying
\eqref{eq:3596-lwe-samples}.  The open question is whether, for standard
parameter families for which LWE has worst-case lattice-hardness
guarantees, there exists a bounded-error quantum algorithm whose running
time is polynomial in $n$ and $\log q$.

\subsection*{Status}

Unsolved

\subsection*{Source}

Contributor: unknown.  The statement is a contributor-formulated open
question synthesizing the quantum-complexity question studied throughout
the LWE literature; it should not be attributed verbatim to Regev, Chen, or
any of the other cited authors.

\subsection*{Progress}

\begin{itemize}
  \item Regev introduced LWE and established its central worst-case hardness
connection: in appropriate parameter regimes, an efficient algorithm for LWE
can be converted, via a quantum reduction, into an efficient quantum
algorithm for approximating worst-case GapSVP and SIVP within polynomial
approximation factors.  The reduction is a worst-case-to-average-case
hardness reduction using quantum computation, not a quantum algorithm for
solving LWE
\sourcecite{ref:3596-regev}{Reg09}.

  \item Brakerski, Kirshanova, Stehl\'e, and Wen introduced the Extrapolated
Dihedral Coset Problem (EDCP) and established quantum polynomial-time
reductions relating LWE to suitable EDCP variants; a polynomial-time
quantum algorithm for EDCP in the arising parameter regime would yield a
polynomial-time quantum algorithm for LWE
\sourcecite{ref:3596-bksw}{BKSW18}.

  \item Chen, Liu, and Zhandry gave polynomial-time quantum algorithms for
variants of average-case lattice problems, including LWE-like inputs in
which the algorithm receives quantum states rather than ordinary classical
LWE samples; the parameter regimes and input models considered are not
known to inherit the standard worst-case hardness guarantees
\sourcecite{ref:3596-clz}{CLZ22}.

  \item Chen proposed a polynomial-time quantum algorithm for LWE with certain
polynomial modulus-to-noise ratios, based on complex Gaussian states and
windowed quantum Fourier transforms; the claim was subsequently withdrawn
after a bug in Step~9 of the algorithm was identified
\sourcecite{ref:3596-chen24}{Che24}.

  \item Chen, Hu, Liu, Luo, and Tu studied LWE variants with quantum amplitudes,
giving a subexponential-time algorithm for certain Gaussian-amplitude LWE
states with known phase and polynomial-time algorithms for quadratic-phase
variants, together with polynomial-time quantum reductions from standard
LWE and worst-case GapSVP to Gaussian-amplitude LWE states carrying a small
unknown phase; the unknown phase is identified as the principal obstruction
to converting these into an efficient algorithm for standard LWE
\sourcecite{ref:3596-chllt}{CHLLT25}.

  \item Bai, Jangir, Kirshanova, Ngo, and Youmans gave a
quasi-polynomial-time quantum algorithm for EDCP over power-of-two moduli,
using a quasi-polynomial number of EDCP states; known reductions from
standard LWE provide only polynomially many relevant states, so this does
not currently imply a quasi-polynomial-time algorithm for standard LWE
\sourcecite{ref:3596-bjkn}{BJKNY25}.
\end{itemize}

\subsection*{References}

\begin{enumerate}[label={},leftmargin=4.8em,labelsep=0.5em]
  \footnotesize
  \item[\textup{[Reg09]}]\label{ref:3596-regev}
  Oded Regev, ``On Lattices, Learning with Errors, Random Linear Codes, and
Cryptography,'' \emph{Journal of the ACM} \textbf{56}, Article 34 (2009).
\href{https://doi.org/10.1145/1568318.1568324}{doi:10.1145/1568318.1568324}.

  \item[\textup{[BKSW18]}]\label{ref:3596-bksw}
  Zvika Brakerski, Elena Kirshanova, Damien Stehl\'e, and Weiqiang Wen,
``Learning With Errors and Extrapolated Dihedral Cosets,'' in \emph{Public-Key
Cryptography -- PKC 2018}, 702--727 (2018).
\href{https://doi.org/10.1007/978-3-319-76581-5_24}{doi:10.1007/978-3-319-76581-5_24};
\href{https://arxiv.org/abs/1710.08223}{arXiv:1710.08223}.

  \item[\textup{[CLZ22]}]\label{ref:3596-clz}
  Yilei Chen, Qipeng Liu, and Mark Zhandry, ``Quantum Algorithms for
Variants of Average-Case Lattice Problems via Filtering,'' in \emph{Advances
in Cryptology -- EUROCRYPT 2022}, 372--401 (2022).
\href{https://doi.org/10.1007/978-3-031-07082-2_14}{doi:10.1007/978-3-031-07082-2_14};
\href{https://arxiv.org/abs/2108.11015}{arXiv:2108.11015}.

  \item[\textup{[Che24]}]\label{ref:3596-chen24}
  Yilei Chen, ``Quantum Algorithms for Lattice Problems,'' IACR Cryptology
ePrint Archive, Paper 2024/555 (2024).
\href{https://eprint.iacr.org/2024/555}{ePrint 2024/555}.

  \item[\textup{[CHLLT25]}]\label{ref:3596-chllt}
  Yilei Chen, Zihan Hu, Qipeng Liu, Han Luo, and Yaxin Tu, ``LWE with
Quantum Amplitudes: Algorithm, Hardness, and Oblivious Sampling,'' in
\emph{Advances in Cryptology -- CRYPTO 2025}, 513--544 (2025).
\href{https://doi.org/10.1007/978-3-032-01878-6_17}{doi:10.1007/978-3-032-01878-6_17};
\href{https://arxiv.org/abs/2310.00644}{arXiv:2310.00644}.

  \item[\textup{[BJKNY25]}]\label{ref:3596-bjkn}
  Shi Bai, Hansraj Jangir, Elena Kirshanova, Tran Ngo, and William Youmans,
``A Quasi-polynomial Time Algorithm for the Extrapolated Dihedral Coset
Problem over Power-of-Two Moduli,'' in \emph{Advances in Cryptology -- CRYPTO
2025}, 416--448 (2025).
\href{https://doi.org/10.1007/978-3-032-01878-6_14}{doi:10.1007/978-3-032-01878-6_14}.
\end{enumerate}

\subsection*{Comment}

Several concrete routes remain open: solving the relevant EDCP instances
with only polynomially many quantum states; overcoming the unknown-phase
obstruction in quantum-amplitude formulations of LWE; repairing or
replacing the complex-Gaussian and windowed-QFT approach; or finding a
fundamentally different quantum algorithm operating directly on ordinary
classical LWE samples.

\subsection*{Field}

Quantum algorithm

\subsection*{Topic}

Computational complexity and computability

\subsection*{ID}

\texttt{op\_3596f8c955c66d94}
